\section{集合、映射与拓扑空间}
流形的“局部像欧氏空间”首先是拓扑性质。本章中$A\subseteq X$允许$A=X$，$2^X$是$X$的幂集。映射$f:X\to Y$的定义域、陪域属于映射的数据；只有满足对应条件的元素才能代入。
\begin{definition}[像、原像与双射]
对$A\subseteq X$、$B\subseteq Y$，定义
\[
f(A)=\{f(a):a\in A\},\qquad f^{-1}(B)=\{x\in X:f(x)\in B\}.
\]
若$f(x)=f(x')$蕴含$x=x'$，称$f$单射；若$f(X)=Y$，称$f$满射；二者兼有时称双射。符号$f^{-1}(B)$表示集合的原像，并不要求$f$有逆映射。
\end{definition}
\begin{theorem}[原像的集合运算]
原像保持任意并、任意交和补集：
\begin{align*}
f^{-1}\Bigl(\bigcup_\alpha B_\alpha\Bigr)&=\bigcup_\alpha f^{-1}(B_\alpha),\\
f^{-1}\Bigl(\bigcap_\alpha B_\alpha\Bigr)&=\bigcap_\alpha f^{-1}(B_\alpha),\\
f^{-1}(Y\setminus B)&=X\setminus f^{-1}(B).
\end{align*}
\end{theorem}
\begin{proof}
$x$属于第一式左边，当且仅当$f(x)$属于某个$B_\alpha$，当且仅当$x$属于某个$f^{-1}(B_\alpha)$，即属于右边。将“某个”改为“每个”得到第二式。第三式两侧都表示满足$f(x)\notin B$的$x\in X$。
\end{proof}
\begin{definition}[拓扑空间]
集合$X$上的拓扑是集合族$\mathcal T\subseteq2^X$，满足：$\varnothing,X\in\mathcal T$；$\mathcal T$对任意并封闭；$\mathcal T$对有限交封闭。$(X,\mathcal T)$称为拓扑空间，$\mathcal T$中的元素称为开集。若$X\setminus F$开，则称$F$闭。
\end{definition}
有限交不能换成无限交：在通常拓扑下，$\bigcap_{n\geq1}(-1/n,1/n)=\{0\}$不是开集。开、闭也不是互补的两个类别；$\varnothing,X$均既开又闭，$[0,1)$在$\R$中则既不开也不闭。
\begin{example}[三种基本拓扑]
离散拓扑是$2^X$，平凡拓扑又称凝聚拓扑，是$\{\varnothing,X\}$。对$\R^n$，通常拓扑由所有满足下述条件的$U$组成：每个$x\in U$都有$r>0$，使开球$B(x,r)\subseteq U$。这确为拓扑：空集条件无须验证，全空间含任意小球；并集中一点落在某个成员中，可沿用该成员的球；有限个开集之交中的点，可取各半径的最小值作为半径。
\end{example}
\begin{theorem}[闭集运算]
任意闭集的交闭，有限闭集的并闭，$\varnothing,X$闭。
\end{theorem}
\begin{proof}
De Morgan公式给出$X\setminus\bigcap F_\alpha=\bigcup(X\setminus F_\alpha)$和$X\setminus\bigcup_{j=1}^mF_j=\bigcap_{j=1}^m(X\setminus F_j)$。右侧分别由任意并、有限交的拓扑公理为开集，故对应左侧的补集闭。最后两个结论来自$\varnothing,X$互为补集。
\end{proof}

\section{邻域、闭包、内部与边界}
\begin{definition}[邻域与集合的三个基本操作]
$N\subseteq X$称为$x$的邻域，若存在开集$U$使$x\in U\subseteq N$。邻域本身不必开。定义
\[
\operatorname{Int}_X A=\bigcup_{U\text{开},\,U\subseteq A}U,\qquad
\overline A^{\,X}=\bigcap_{F\text{闭},\,A\subseteq F}F,\qquad
\partial_XA=\overline A^{\,X}\setminus\operatorname{Int}_XA.
\]
无歧义时省略环境空间$X$。内部是包含于$A$的最大开集，闭包是包含$A$的最小闭集；这些结论分别由开集并、闭集交的性质直接得到。
\end{definition}
\begin{theorem}[邻域刻画与边界]
集合$A$开，当且仅当它是其中每一点的邻域。并且
\[
x\in\overline A\iff\text{$x$的每个开邻域都与$A$相交},\qquad
\partial A=\overline A\cap\overline{X\setminus A}.
\]
因此边界是闭集；$A$开当且仅当$A=\operatorname{Int}A$，$A$闭当且仅当$A=\overline A$。
\end{theorem}
\begin{proof}
若$A$开，取$U=A$即可。反之，对每个$x\in A$取开集$U_x$满足$x\in U_x\subseteq A$，则$A=\bigcup_{x\in A}U_x$开。
若$x\notin\overline A$，开集$X\setminus\overline A$包含$x$且不交$A$。反之，若开邻域$U$不交$A$，则闭集$X\setminus U$包含$A$，从而$\overline A\subseteq X\setminus U$，所以$x\notin\overline A$。这证明闭包刻画。
点$x$不属于$\overline{X\setminus A}$，恰好表示有开邻域$U\subseteq A$，也就是$x\in\operatorname{Int}A$。取补集并与$\overline A$相交，便得边界公式。两闭集之交闭。最后两项由内部和闭包的最大、最小性质及其开闭性得到。
\end{proof}
\begin{example}[环境空间改变，内部和边界也改变]
令$A=[0,1)\subseteq\R$。在$\R$中，内部为$(0,1)$、闭包为$[0,1]$、边界为$\{0,1\}$；在空间$A$自身中，$A$是全空间，故其内部和闭包均为$A$，边界为空。这提示我们必须先说明拓扑所在的空间。
\end{example}

\section{子空间拓扑与相对开闭集}
\begin{definition}[子空间拓扑]
若$(X,\mathcal T)$为拓扑空间且$A\subseteq X$，定义
\[
\mathcal T|_A=\{A\cap U:U\in\mathcal T\}.
\]
它称为$A$的子空间拓扑，也称由$X$诱导的拓扑。以后把欧氏空间的子集看作拓扑空间时，除非另行说明，都使用此拓扑。
\end{definition}
\begin{theorem}[子空间拓扑的合法性与传递性]\label{thm:subspace}
$\mathcal T|_A$是拓扑。若$B\subseteq A\subseteq X$，先从$X$诱导到$A$再到$B$，与直接从$X$诱导到$B$所得拓扑相同。并且$F\subseteq A$在$A$中闭，当且仅当$F=A\cap C$，其中$C$在$X$中闭。
\end{theorem}
\begin{proof}
$\varnothing=A\cap\varnothing$、$A=A\cap X$在该集合族中。任意并满足$\bigcup_\alpha(A\cap U_\alpha)=A\cap\bigcup_\alpha U_\alpha$；有限交满足$\bigcap_{j=1}^m(A\cap U_j)=A\cap\bigcap_{j=1}^mU_j$。括号后的并或交在$X$中开，故三条公理成立。
在$B$中由两次诱导得到的开集形如$B\cap(A\cap U)=B\cap U$；这正是一次诱导得到的全部开集。
最后，若$A\setminus F=A\cap U$且$U$在$X$中开，则$F=A\cap(X\setminus U)$，取$C=X\setminus U$即可。反之，若$F=A\cap C$且$C$闭，则$A\setminus F=A\cap(X\setminus C)$相对开，因此$F$相对闭。
\end{proof}
\begin{theorem}[子空间闭包与开子空间]
若$B\subseteq A\subseteq X$，则$\overline B^{\,A}=A\cap\overline B^{\,X}$。若$A$在$X$中开，则$V\subseteq A$在$A$中开当且仅当它在$X$中开。
\end{theorem}
\begin{proof}
对$x\in A$，它在$A$中的任意开邻域形如$A\cap U$，其中$U$在$X$中开且$x\in U$。由于$B\subseteq A$，$(A\cap U)\cap B=U\cap B$。利用闭包的邻域刻画，$x$属于左侧当且仅当属于右侧。
若$V=A\cap U$，$A,U$都在$X$中开，则$V$在$X$中开；反之若$V\subseteq A$在$X$中开，则$V=A\cap V$相对开。
\end{proof}
\begin{example}[相对开不等于环境中开]
$[0,\tfrac12)$在$A=[0,1]$中相对开，因为它等于$A\cap(-1,\tfrac12)$，但它在$\R$中不开放。在圆周$S^1\subseteq\R^2$上，一段不含端点的短圆弧相对开，却不是平面开集，因为任何以弧上点为中心的平面开球都含不在圆周上的点。
\end{example}
\begin{figure}[htbp]
\centering
\begin{tikzpicture}[scale=1.5,>=Stealth]
\draw[->,gray] (-1.6,0)--(2.4,0) node[right] {$x$};
\draw[->,gray] (0,-1.35)--(0,1.4) node[above] {$y$};
\fill[blue!12] (1,0) circle (.8);
\draw[blue,dashed] (1,0) circle (.8);
\draw[thick] (0,0) circle (1);
\draw[line width=2pt,orange!85!black] (47.16:1) arc (47.16:-47.16:1);
\fill[white] (47.16:1) circle (.035);
\fill[white] (-47.16:1) circle (.035);
\draw[orange!85!black] (47.16:1) circle (.035);
\draw[orange!85!black] (-47.16:1) circle (.035);
\node at (-.9,.95) {$A=S^1$};
\node[blue] at (1.65,.85) {$U$};
\node[orange!85!black] at (1.6,-1) {$A\cap U$};
\end{tikzpicture}

\caption{参照梁灿彬、周彬上册第二版印刷页7（PDF第23页）图1-4重新绘制：圆周中的相对开弧是平面开集与圆周的交。端点不属于开圆盘，也不属于该开弧。}
\end{figure}

\section{连续映射、同胚与拓扑基}
\begin{definition}[连续与同胚]
映射$f:X\to Y$连续，是指$Y$中每个开集的原像在$X$中开。若$f$是双射且$f,f^{-1}$都连续，则称$f$为同胚。两空间之间存在同胚时，称两空间同胚。
\end{definition}
\begin{theorem}[连续性的等价刻画]\label{thm:continuity}
对$f:X\to Y$，下列条件等价：开集原像开；闭集原像闭；对任意$x\in X$和$f(x)$的开邻域$V$，有$x$的开邻域$U$满足$f(U)\subseteq V$。
\end{theorem}
\begin{proof}
由$f^{-1}(Y\setminus F)=X\setminus f^{-1}(F)$，开集原像开与闭集原像闭等价。若开集原像开，取$U=f^{-1}(V)$即得邻域条件。反之，给定开集$V\subseteq Y$，对每个$x\in f^{-1}(V)$取邻域条件中的$U_x$，则$U_x\subseteq f^{-1}(V)$，所以$f^{-1}(V)=\bigcup_{x\in f^{-1}(V)}U_x$开。
\end{proof}
\begin{theorem}[复合、限制与取值于子空间]\label{thm:restrict}
连续映射的复合连续。若$f:X\to Y$连续，则其限制$f|_A:A\to Y$连续。若$f(X)\subseteq B\subseteq Y$，则$f:X\to B$连续当且仅当$f:X\to Y$连续，这里$A,B$均取子空间拓扑。
\end{theorem}
\begin{proof}
对连续映射$f:X\to Y$和$g:Y\to Z$及开集$W\subseteq Z$，$(g\circ f)^{-1}(W)=f^{-1}(g^{-1}(W))$是开集。限制映射满足$(f|_A)^{-1}(V)=A\cap f^{-1}(V)$，故连续。
若$f(X)\subseteq B$，则对$B$中开集$B\cap V$，有$f^{-1}(B\cap V)=f^{-1}(V)$，这证明$Y$值连续蕴含$B$值连续。反过来，对$Y$中开集$V$，$f^{-1}(V)=f^{-1}(B\cap V)$，若$B$值映射连续，右侧开。
\end{proof}
\begin{definition}[拓扑基与可数性]
一族开集$\mathcal B$称为拓扑基，若每个开集都是其中成员的并；等价地，对每个开集$U$和$x\in U$，存在$B\in\mathcal B$使$x\in B\subseteq U$。存在可数拓扑基的空间称为第二可数空间。若对每个点$x$，有一列开邻域使$x$的任意开邻域都包含其中某个，则称空间第一可数。
\end{definition}
\begin{theorem}[欧氏空间的可数基及其遗传性]\label{thm:countable-base}
$\R^n$中有理中心、正有理半径的开球构成可数基。第二可数空间的任意子空间第二可数；Hausdorff空间的任意子空间Hausdorff。
\end{theorem}
\begin{proof}
给定$x\in U$且$U$开，取$\varepsilon>0$使$B(x,\varepsilon)\subseteq U$。由有理数的稠密性，取$q\in\mathbb Q^n$使$|q-x|<\varepsilon/4$，再取有理数$r$使$\varepsilon/4<r<\varepsilon/2$。则$x\in B(q,r)$，而$|y-x|\leq|y-q|+|q-x|<3\varepsilon/4$说明此球包含于$U$。这些球由可数集$\mathbb Q^n\times\mathbb Q_{>0}$编号。
若$\{B_j\}$是$X$的可数基，则$\{A\cap B_j\}$是子空间$A$的可数基，因为任意相对开集$A\cap U$可按$U$的基展开。若$x,y\in A$不同，在$X$中用不交开集分离它们，再分别与$A$相交，即在$A$中分离它们。
\end{proof}
\begin{definition}[乘积与商拓扑]
$X\times Y$上的乘积拓扑以$U\times V$为基，其中$U,V$分别开。若$q:X\to Q$为满射，规定$W\subseteq Q$开当且仅当$q^{-1}(W)$在$X$中开，所得拓扑称为商拓扑。
\end{definition}
\begin{theorem}[乘积与商的映射判别]\label{thm:product-quotient}
乘积拓扑及商拓扑都是拓扑。$h:Z\to X\times Y$连续当且仅当其两个分量连续。对商映射$q:X\to Q$，映射$k:Q\to Z$连续当且仅当$k\circ q$连续。
\end{theorem}
\begin{proof}
开矩形覆盖乘积，两个开矩形的交仍是开矩形，所以它们的任意并满足拓扑公理。商拓扑的公理由原像保持并、交和补集直接验证。
乘积投影连续，因为$\pi_X^{-1}(U)=U\times Y$开；故$h$连续则分量连续。若分量连续，则$h^{-1}(U\times V)=h_X^{-1}(U)\cap h_Y^{-1}(V)$开，基中开集的原像开进一步推出所有开集原像开。
若$k$连续，由$q$连续知复合连续。反之，对$Z$中开集$W$，$q^{-1}(k^{-1}(W))=(k\circ q)^{-1}(W)$开，由商拓扑定义得$k^{-1}(W)$开。
\end{proof}

\section{连通性与道路连通性}
\begin{definition}[连通与道路连通]
空间$X$连通，是指不能写为两个非空、不交开集的并。等价地，既开又闭的子集只有$\varnothing,X$。道路是连续映射$\gamma:[0,1]\to X$；若任意两点都是某条道路的端点，则$X$道路连通。子集的连通性一律按子空间拓扑判断。
\end{definition}
等价性来自如下事实：若$X=U\sqcup V$且二者开，则$U$的补集$V$也开，所以$U$闭；反过来，一个非平凡既开又闭集及其补集给出分离。
\begin{theorem}[区间、连续像与道路]
实数区间连通。连通空间的连续像连通。道路连通空间连通。
\end{theorem}
\begin{proof}
设区间$I$被相对开集$U,V$分离。各取一点，交换命名后可设$u\in U,v\in V,u<v$。令$c=\sup(U\cap[u,v])$，则$c\in[u,v]\subseteq I$。若$c\in U$，相对开放性及$v\notin U$说明$c<v$，并在$c$右侧还能找到$U\cap[u,v]$中的点，与上界矛盾。若$c\in V$，因$u\notin V$有$c>u$，相对开放性给出$c$左侧某区间不含$U$中的点，与$c$为上确界矛盾。所以分离不存在。
若$f(X)$被两个非空相对开集分离，它们在$X$中的原像也非空、不交、开且覆盖$X$，与连通性矛盾。
最后，若道路连通的$X$存在分离，分别取两侧的点并用道路连接。$[0,1]$连通，其连续像不能同时与分离两侧相交，矛盾。
\end{proof}

\section{紧致性、Hausdorff性与基本定理}
\begin{definition}[开覆盖、紧致与Hausdorff]
集合族$\{U_\alpha\}$是$A\subseteq X$的开覆盖，若各$U_\alpha$在$X$中开且$A\subseteq\bigcup U_\alpha$。若每个开覆盖都有有限子覆盖，称$A$紧致。此定义等价于$A$作为子空间紧致，因为相对开集均形如$A\cap U$。若不同两点均有互不相交的开邻域，称空间Hausdorff。
\end{definition}
\begin{theorem}[紧致性的三条工具]\label{thm:compact-tools}
紧致空间的闭子集紧致；连续映射把紧致集映成紧致集；Hausdorff空间中的紧致子集闭。
\end{theorem}
\begin{proof}
若$F$是紧致空间$X$的闭子集，将$F$的相对开覆盖写为$\{F\cap U_\alpha\}$，再用$X\setminus F$补成$X$的开覆盖。取有限子覆盖并去掉$X\setminus F$，其余对应成员覆盖$F$。
若$K$紧致，$\{V_\alpha\}$是$f(K)$的相对开覆盖，各原像覆盖$K$；由限制及子空间连续性，它们在$K$中开。有限个原像覆盖$K$意味着相应的有限个$V_\alpha$覆盖$f(K)$。
若$K$是Hausdorff空间$Y$中的非空紧致集且$y\notin K$，对每个$z\in K$取互不相交的开邻域$U_z\ni z$、$V_z\ni y$。有限个$U_{z_1},\ldots,U_{z_m}$覆盖$K$。则$\bigcap_{j=1}^mV_{z_j}$是$y$的开邻域，且与$K$不交，所以$Y\setminus K$开。空集的情形直接成立。
\end{proof}
\begin{theorem}[紧致到Hausdorff的连续双射]\label{thm:compact-homeomorphism}
若$X$紧致、$Y$为Hausdorff空间，且$f:X\to Y$是连续双射，则$f$为同胚。
\end{theorem}
\begin{proof}
对$X$中闭集$F$，上一定理依次给出$F$紧致、$f(F)$紧致、$f(F)$在$Y$中闭。令$h=f^{-1}$，则$h^{-1}(F)=f(F)$闭。由定理\ref{thm:continuity}，$h$连续，故$f$是同胚。
\end{proof}
\begin{theorem}[Heine--Borel定理]\label{thm:heine-borel}
欧氏空间$\R^n$中的子集紧致，当且仅当它闭且有界。
\end{theorem}
\begin{proof}
紧致集$K$在Hausdorff空间$\R^n$中闭。开球族$\{B(0,m):m\in\mathbb N,m\geq1\}$覆盖$K$，有限子覆盖中的最大半径说明$K$有界。
反之，先证明闭立方体$C=[-M,M]^n$紧致，其中$M>0$。若某开覆盖没有有限子覆盖，将$C$沿每个坐标中点分为$2^n$个闭小立方体，至少一个仍不能被有限覆盖，否则合并这些有限覆盖就能覆盖$C$。重复选取，得到嵌套闭立方体$C_j$，边长趋于零。对每个坐标，它们的投影是嵌套闭区间$[a_j,b_j]$；令$a=\sup_j a_j$，嵌套性保证$a_j\leq a\leq b_j$，故各坐标给出共同点$x\in\bigcap_jC_j$。取覆盖中包含$x$的某开集$U$，有$\varepsilon>0$使$B(x,\varepsilon)\subseteq U$。当$C_j$的直径小于$\varepsilon$时，由$x\in C_j$知$C_j\subseteq U$，与$C_j$不能有限覆盖矛盾。
若$K$闭且有界，选择$M>0$使$K\subseteq C$。它是紧致空间$C$的闭子集，故紧致。
\end{proof}
\begin{theorem}[有限乘积紧致]
若$X,Y$紧致，则$X\times Y$在乘积拓扑下紧致。
\end{theorem}
\begin{proof}
空因子的情形乘积为空，结论成立。给定乘积的开覆盖$\{W_\alpha\}$，对每个$(x,y)$选包含它的开矩形$U_{x,y}\times V_{x,y}\subseteq W_{\alpha(x,y)}$。固定$x$，$\{V_{x,y}:y\in Y\}$覆盖$Y$，故可选有限个$y_1,\ldots,y_m$覆盖它。置$U_x=\bigcap_{j=1}^mU_{x,y_j}$，则$U_x\times Y$被上述有限个$W_{\alpha(x,y_j)}$覆盖。再由$X$紧致，有限个$U_x$覆盖$X$。把这些有限列表合并即得整个乘积的有限子覆盖。归纳得到任意有限乘积的结论。
\end{proof}
\begin{example}[连续双射的限制]
$q:[0,2\pi)\to S^1$，$q(t)=(\cos t,\sin t)$是连续双射，却非同胚。$[0,\varepsilon)$是定义域的相对开集，而它的像在圆周中不是开集：包含$(1,0)$的每个圆周开邻域还含辐角接近$2\pi$的点。这里定义域不紧致，所以不满足定理\ref{thm:compact-homeomorphism}的条件。
\end{example}
